\chapter{Energi: Bentuk dan Kekekalannya}\label{ch:g10:energy-conservation}

Mobil yang melaju, telepon yang terisi, oven yang panas dan palu yang
terangkat berbagi satu sifat: masing-masing dapat membuat sesuatu
terjadi. Fisika mengukur kesanggupan itu dengan satu bilangan ---
\omterm{def:g10:energy-conservation:energy}{energi} --- dan menutup pembukuannya dengan satu aturan besi: jumlah
seluruhnya tidak pernah berubah, ia hanya berpindah dan berganti
busana. Bab ini membuka buku besarnya: bentuk yang diambil \omterm{def:g10:energy-conservation:energy}{energi},
rantai yang ditempuhnya, dan apa yang dibeli satu \omterm{def:g10:energy-conservation:kwh}{kilowatt-jam}.

\section{Energi: mata uang bersama}

\begin{definition}[Energi dan joule]\label{def:g10:energy-conservation:energy}
\emph{Energi}\index{energi} mengukur kesanggupan sebuah sistem untuk
menghasilkan perubahan: menggerakkan materi, mengangkatnya,
memanaskannya, menyalakannya. \omterm{def:g10:orders-of-magnitude:unit}{Satuannya} adalah
\emph{joule}\index{joule} (\unit{J}). Tersimpan, berpindah, berubah
bentuk --- setiap cabang fisika berhitung dengan satu mata uang ini.
\end{definition}

\begin{definition}[Energi kinetik]\label{def:g10:energy-conservation:kinetic}
Sebuah benda bermassa $m$ (dalam \unit{kg}) yang bergerak dengan laju
$v$ (dalam \unit{m/s}) membawa \emph{energi kinetik}\index{energi
kinetik} $E_k = \tfrac12\, m v^2$: menggandakan massanya menggandakan
\omterm{def:g10:energy-conservation:energy}{energinya}, sedangkan menggandakan lajunya melipatempatkannya.
\end{definition}

\begin{example}[Mobil lawan pejalan kaki]\label{ex:g10:energy-conservation:kinetic}
Mobil \qty{1300}{kg} pada \qty{100}{km/h} $= \qty{27.8}{m/s}$ membawa
$E_k = \tfrac12 \times 1300 \times 27.8^2 \approx \qty{5.0e5}{J}$;
seorang pejalan kaki \qty{70}{kg} pada \qty{1.5}{m/s} membawa
$\tfrac12 \times 70 \times 1.5^2 \approx \qty{79}{J}$ --- enam ribu
kali lebih kecil. Laju itu mahal: kuadratnyalah yang mengurusnya.
\end{example}

\begin{definition}[Energi potensial gravitasi]\label{def:g10:energy-conservation:potential}
Sebuah benda bermassa $m$ pada ketinggian $h$ (dalam \unit{m}) di atas
suatu aras acuan yang dipilih menyimpan \emph{energi potensial
gravitasi}\index{energi potensial!gravitasi} $E_p = m g h$, dengan
$g = \qty{9.81}{N/kg}$ --- berlaku di dekat permukaan, tempat $g$ boleh
dianggap tetap. Yang berarti hanyalah \emph{selisih} ketinggiannya:
aras tempat $E_p = 0$ boleh dipilih dengan bebas.
\end{definition}

\begin{example}[Menara air]\label{ex:g10:energy-conservation:potential}
Tiap \omterm{def:g10:orders-of-magnitude:unit}{meter} kubik (\qty{1000}{kg}) air yang tersimpan \qty{30}{m} di
atas sebuah menara air memuat
$1000 \times 9.81 \times 30 \approx \qty{2.9e5}{J}$ --- kira-kira enam
persepuluh \omterm{def:g10:energy-conservation:kinetic}{energi kinetik} mobil di atas tadi. Kota menimbun \omterm{def:g10:energy-conservation:energy}{energi} di
atas kepala.
\end{example}

\begin{definition}[Bentuk yang lain]\label{def:g10:energy-conservation:forms}
\omterm{def:g10:energy-conservation:energy}{Energi} juga bersembunyi dalam busana yang kurang kasatmata:
\emph{energi termal}\index{energi termal}, yaitu keriuhan tak beraturan
molekul sebuah benda (makin panas, makin riuh); \emph{energi
kimia}\index{energi kimia}, yang tersimpan pada ikatan molekul ---
makanan, kayu, bensin, baterai; \emph{energi listrik}\index{energi
listrik}, yang dibawa arus dari pembangkit sampai stopkontak;
\emph{energi radiasi}\index{energi radiasi}, yang dibawa cahaya bahkan
melintasi ruang kosong --- begitulah \omterm{def:g10:energy-conservation:energy}{energi} Matahari sampai kepada
kita; \emph{energi nuklir}\index{energi nuklir}, yang tersimpan di
dalam inti atom dan dilepaskan di dalam bintang
(\cref{ch:g12:nuclear-energy}).
\end{definition}

\section{Perpindahan dan rantai pengubahan}

\begin{definition}[Perpindahan, pengubahan, rantai]\label{def:g10:energy-conservation:transfer}
\emph{Perpindahan energi}\index{perpindahan energi} memindahkan \omterm{def:g10:energy-conservation:energy}{energi}
antarsistem --- antar\emph{penampung}\index{penampung (energi)};
\emph{pengubahan energi}\index{pengubahan energi} mengubah bentuknya.
Alat yang melakukan salah satunya (turbin, motor, lampu, otot) disebut
\emph{pengubah}\index{pengubah}. \emph{Rantai
pengubahan}\index{rantai pengubahan} menggambarkan perjalanannya: kotak
untuk penampung dan bentuknya, anak panah berlabel untuk
perpindahannya.
\end{definition}

\begin{omfigure}
\begin{tikzpicture}[font=\small,
  box/.style={draw, rounded corners, fill=omDef!10, minimum height=10mm, align=center},
  lab/.style={font=\footnotesize, align=center}]
  \node[box] (res)  at (0,0)   {air waduk\\(potensial)};
  \node[box] (tur)  at (4.1,0) {turbin berputar\\(kinetik)};
  \node[box] (grid) at (8.2,0) {rumah pada jaringan\\(listrik)};
  \draw[-Stealth, thick, omThm] (res) -- (tur)
    node[midway, above, lab] {jatuh melalui\\pipa pesat};
  \draw[-Stealth, thick, omThm] (tur) -- (grid)
    node[midway, above, lab] {alternator};
  \draw[-Stealth, omExo] (tur.south) -- ++(0,-0.7)
    node[below, lab] {kalor (gesekan,\\pusaran)};
  \draw[-Stealth, omExo] (grid.south) -- ++(0,-0.7)
    node[below, lab] {kalor di kawatnya};
\end{tikzpicture}

{\small \omterm{def:g10:energy-conservation:transfer}{Rantai pengubahan} sebuah bendungan pembangkit listrik: kotaknya
\omterm{def:g10:energy-conservation:transfer}{penampung}, anak panahnya perpindahan --- setiap anak panah yang
sebenarnya membocorkan kalor.}
\end{omfigure}

\begin{method}[Menggambar rantai pengubahan]\label{met:g10:energy-conservation:chain}
\begin{enumerate}
  \item Kenali \omterm{def:g10:energy-conservation:transfer}{penampung} dan bentuk awalnya (apa yang dikonsumsi?).
  \item Ikuti \omterm{def:g10:energy-conservation:energy}{energinya} dari \omterm{def:g10:energy-conservation:transfer}{pengubah} ke \omterm{def:g10:energy-conservation:transfer}{pengubah}: satu anak panah tiap
        perpindahan, dan bentuknya disebutkan di tiap kotak.
  \item Tambahkan anak panah kalor pada tiap \omterm{def:g10:energy-conservation:transfer}{pengubah} yang sebenarnya
        --- selalu ada \omterm{def:g10:energy-conservation:energy}{energi} yang bocor sebagai kalor (\omterm{def:g10:inertia:inventory}{gesekan}, kawat,
        gas buang).
\end{enumerate}
\end{method}

Sebuah bandul menjalankan rantai yang paling sederhana dari semuanya
--- seluruhnya potensial di puncak ayunannya, seluruhnya kinetik ketika
melewati titik terendahnya (dengan mengambil $E_p = 0$ di sana):

\begin{omfigure}
\begin{tikzpicture}[font=\small]
  \draw[-Stealth] (0,0) -- (0,2.5) node[above] {$E$}; \draw (0,0) -- (2.8,0);
  \draw[dashed, omExo] (0,2.0) -- (2.8,2.0) node[right] {jumlah};
  \fill[omDef!60] (0.5,0) rectangle (1.1,2.0); \draw[omThm, very thick] (1.7,0) -- (2.3,0);
  \node at (0.8,-0.4) {$E_p$}; \node at (2.0,-0.4) {$E_k$};
  \node at (1.4,-1.0) {puncak ayunan};
\end{tikzpicture}\qquad
\begin{tikzpicture}[font=\small]
  \draw[-Stealth] (0,0) -- (0,2.5) node[above] {$E$}; \draw (0,0) -- (2.8,0);
  \draw[dashed, omExo] (0,2.0) -- (2.8,2.0) node[right] {jumlah};
  \draw[omDef, very thick] (0.5,0) -- (1.1,0); \fill[omThm!60] (1.7,0) rectangle (2.3,2.0);
  \node at (0.8,-0.4) {$E_p$}; \node at (2.0,-0.4) {$E_k$};
  \node at (1.4,-1.0) {titik terendah};
\end{tikzpicture}

{\small Anggaran bandul itu pada dua kedudukan: $E_p$ dan $E_k$
bertukar sementara jumlahnya (garis putus-putus) tinggal tetap.}
\end{omfigure}

\section{Kekekalan energi}

\begin{theorem}[Kekekalan energi]\label{thm:g10:energy-conservation:conservation}
\index{kekekalan energi}
\omterm{def:g10:energy-conservation:energy}{Energi} seluruhnya pada sistem terasing --- yaitu sistem yang tidak
menukarkan apa pun dengan luarnya --- bernilai tetap. \omterm{def:g10:energy-conservation:energy}{Energi} tidak
pernah diciptakan dan tidak pernah dimusnahkan: ia hanya berganti
bentuk dan tempat, dan apa yang hilang dari satu \omterm{def:g10:energy-conservation:transfer}{penampung} diperoleh
\omterm{def:g10:energy-conservation:transfer}{penampung} yang lain, \omterm{def:g10:energy-conservation:energy}{joule} demi \omterm{def:g10:energy-conservation:energy}{joule}.
\end{theorem}
\admitted

\begin{remark}\label{rem:g10:energy-conservation:admitted}
Tidak pernah ada percobaan yang menangkap asas ini gagal; di sini kita
menerimanya sebagai diketahui lalu memakainya secara berangka.
\cref{ch:g12:work-and-energy} menurunkan bagian mekanisnya
($E_k + E_p$) dari hukum Newton; pembukuan yang lengkap menunggu jilid
perguruan tinggi.
\end{remark}

\begin{example}[Jatuh bebas, dihargai dalam joule]\label{ex:g10:energy-conservation:fall}
Sebuah bola \qty{2.0}{kg} jatuh dari \qty{10}{m} (hambatan udaranya
diabaikan): ia berangkat dengan
$E_p = 2.0 \times 9.81 \times 10 \approx \qty{196}{J}$, mendarat dengan
$E_k = \qty{196}{J}$, sehingga $v = \sqrt{2 \times 196/2.0} =
\qty{14}{m/s} \approx \qty{50}{km/h}$. Tanpa \omterm{def:g10:inertia:force}{gaya}, tanpa pencatat
waktu: kekekalannya sendirian sudah menghargai tumbukan itu.
\end{example}

\section{Efisiensi}

\begin{definition}[Efisiensi]\label{def:g10:energy-conservation:efficiency}
Sebuah \omterm{def:g10:energy-conservation:transfer}{pengubah} mengonsumsi $E_{\text{terpakai}}$ dan menyerahkan
bentuk yang diinginkan $E_{\text{berguna}}$;
\emph{efisiensi}\index{efisiensi}nya adalah
\[
\eta = \frac{E_{\text{berguna}}}{E_{\text{terpakai}}} < 1 .
\]
Kekekalan melarang $\eta > 1$, dan \omterm{def:g10:energy-conservation:transfer}{pengubah} yang sebenarnya tidak
pernah mencapai $1$: \omterm{def:g10:energy-conservation:energy}{energi} yang hilang itu tidak dimusnahkan --- ia
pergi sebagai kalor.
\end{definition}

\begin{example}[Ke mana bensinnya pergi]\label{ex:g10:energy-conservation:engine}
Satu \omterm{def:g10:orders-of-magnitude:unit}{kilogram} bensin menyimpan \qty{45}{MJ} \omterm{def:g10:energy-conservation:forms}{energi kimia}; sebuah mesin
mobil mengubahnya dengan $\eta \approx 0.30$: \qty{13.5}{MJ} menjadi
gerak, \qty{31.5}{MJ} memanaskan gas buangnya, radiatornya dan
jalannya. Dua pertiga tiap tangki menghangatkan atmosfer.
\end{example}

\begin{omfigure}
\begin{tikzpicture}[font=\small]
  \fill[omDef!40] (0,0) rectangle (3,1.5);
  \node[align=center] at (1.5,0.75) {bensin\\\qty{45}{MJ}};
  \fill[omMeth!60] (3,1.05) -- (5.6,1.55) -- (5.6,2.0) -- (3,1.5);
  \fill[omMeth!60] (5.6,1.45) -- (6.3,1.775) -- (5.6,2.1);
  \node[right] at (6.35,1.78) {gerak: \qty{13.5}{MJ} ($30\%$)};
  \fill[omExo!70] (3,0) -- (5.6,-1.0) -- (5.6,0.05) -- (3,1.05);
  \fill[omExo!70] (5.6,-1.1) -- (6.3,-0.475) -- (5.6,0.15);
  \node[right] at (6.35,-0.48) {kalor: \qty{31.5}{MJ} ($70\%$)};
\end{tikzpicture}

{\small Anggaran mesin mobil itu, dengan lebar anak panah sesuai skala:
dari \qty{45}{MJ} yang dikonsumsi, \qty{13.5}{MJ} menggerakkan
mobilnya; sisanya menjadi kalor.}
\end{omfigure}

\begin{remark}[Efisiensi saling dikalikan]\label{rem:g10:energy-conservation:chain-eta}
\omterm{def:g10:energy-conservation:transfer}{Pengubah} yang dirangkai mengalikan \omterm{def:g10:energy-conservation:efficiency}{efisiensinya}: sebuah pembangkit
($0.38$) yang memasok jaringan ($0.94$) yang memasok pengisi \omterm{def:g10:energy-conservation:power}{daya}
($0.85$) menyerahkan $\eta = 0.38 \times 0.94 \times 0.85 \approx 0.30$. Rantai yang panjang adalah pipa yang bocor: tiap anak panah
memungut pajak dari alirannya.
\end{remark}

\section{Daya, watt dan kilowatt-jam}

\begin{definition}[Daya]\label{def:g10:energy-conservation:power}
\emph{Daya}\index{daya} adalah laju mengalirnya \omterm{def:g10:energy-conservation:energy}{energi}, $P = E/t$, yang
diukur dalam \emph{watt}\index{watt} (\unit{W}), dengan
$\qty{1}{W} = \qty{1}{J/s}$. \omterm{def:g10:energy-conservation:energy}{Energi} adalah sebuah jumlah, daya adalah
sebuah aliran --- bak mandi berbanding kerannya.
\end{definition}

\begin{example}[Ketel lawan manusia]\label{ex:g10:energy-conservation:power}
Ketel \qty{2200}{W} yang menyala \qty{150}{s} memakai
$E = 2200 \times 150 = \qty{3.3e5}{J}$. Seorang manusia berjalan dengan
sekitar \qty{10}{MJ} makanan tiap hari, yaitu rata-rata
$\num{1.0e7}/\num{86400} \approx \qty{116}{W}$: kamu menyala seperti
sebuah bohlam tua yang terang.
\end{example}

\begin{definition}[Kilowatt-jam]\label{def:g10:energy-conservation:kwh}
\emph{Kilowatt-jam}\index{kilowatt-jam} (\unit{kW.h}) adalah \omterm{def:g10:energy-conservation:energy}{energi}
dari aliran \qty{1}{kW} yang berlangsung satu jam:
$\qty{1}{kW.h} = 1000 \times 3600\,\unit{J} = \qty{3.6}{MJ}$ --- inilah
\omterm{def:g10:orders-of-magnitude:unit}{satuan} yang dihitung meteran listrik: \omterm{def:g10:energy-conservation:power}{daya} dalam \unit{kW} dikali jam.
\end{definition}

\begin{example}[Membaca tagihan]\label{ex:g10:energy-conservation:bill}
Sebuah tagihan adalah satu pengurangan dan satu perkalian: bacaan
meteran yang baru dikurangi yang lama memberikan \omterm{def:g10:energy-conservation:kwh}{kilowatt-jamnya};
dikalikan harga \omterm{def:g10:orders-of-magnitude:unit}{satuannya} (sekitar $0.25$ euro tiap \unit{kW.h}),
ditambah langganan tetapnya, memberikan jumlah seluruhnya. Pemanas
\qty{2}{kW} yang menyala \qty{2.5}{h} tiap malam memakai \qty{5}{kW.h}
tiap hari: $38$ euro sebulan hanya untuk pemanas itu.
\end{example}

\begin{example}[Orde besar]\label{ex:g10:energy-conservation:orders}
Patut dihafal: mengangkat sebuah apel setinggi satu \omterm{def:g10:orders-of-magnitude:unit}{meter}, sekitar
\qty{1}{J}; baterai telepon, $\qty{10}{W.h} \approx \qty{4e4}{J}$;
semangkuk sereal, \qty{1.5}{MJ}; makananmu untuk sehari, \qty{10}{MJ};
satu \omterm{def:g10:orders-of-magnitude:unit}{kilogram} bensin, \qty{45}{MJ} --- hampir persis sama dengan
listrik harian sebuah rumah tangga ($\approx \qty{12}{kW.h}$); sebuah
sambaran petir, \qty{5}{GJ}: seratus hari rumah tangga, yang
diserahkan dalam sekejap.
\end{example}

\begin{omfigure}
\begin{tikzpicture}[font=\small, y=0.5cm]
  \draw[-Stealth] (0,-0.5) -- (0,10.8) node[above] {$E$ (\unit{J})};
  \foreach \n in {0,2,4,6,8,10}
    {\draw (-0.12,\n) -- (0.12,\n); \node[left, font=\footnotesize] at (-0.15,\n) {$10^{\n}$};}
  \fill[omDef] (0,0) circle (1.6pt) node[right=4pt] {apel diangkat \qty{1}{m}};
  \fill[omDef] (0,4.56) circle (1.6pt) node[right=4pt] {baterai telepon (\qty{10}{W.h})};
  \fill[omDef] (0,6.18) circle (1.6pt) node[right=4pt] {semangkuk sereal};
  \fill[omDef] (0,7) circle (1.6pt) node[right=4pt] {makananmu, sehari (\qty{10}{MJ})};
  \fill[omThm] (0,7.65) circle (1.6pt) node[right=4pt] {\qty{1}{kg} bensin $\approx$ sehari rumah tangga};
  \fill[omThm] (0,9.7) circle (1.6pt) node[right=4pt] {sambaran petir ($\approx\qty{5}{GJ}$)};
\end{tikzpicture}

{\small Sebuah tangga \omterm{def:g10:energy-conservation:energy}{energi}: tiap anak tangganya berfaktor $100$.
Kehidupan sehari-hari terbentang sepuluh \omterm{def:g10:orders-of-magnitude:oom}{orde besar}.}
\end{omfigure}

\section{Latihan}

\begin{exercise}[$\star$]\label{exo:g10:energy-conservation:1}
Sebuah mobil \qty{1200}{kg} melaju \qty{90}{km/h}, yaitu \qty{25}{m/s}:
hitunglah \omterm{def:g10:energy-conservation:kinetic}{energi kinetiknya}. Ulangi pada \qty{130}{km/h}
(\qty{36.1}{m/s}): apa yang dilakukan laju yang $44\%$ lebih besar itu
terhadap $E_k$?
\end{exercise}

\begin{exercise}[$\star$]\label{exo:g10:energy-conservation:2}
Seorang pendaki bermassa \qty{72}{kg} naik dari ketinggian
\qty{1200}{m} ke \qty{1650}{m}. Hitunglah \omterm{def:g10:energy-conservation:energy}{energi} potensial yang
diperolehnya, ubahlah ke \omterm{def:g10:energy-conservation:kwh}{kilowatt-jam}, lalu hargailah pada $0.25$ euro
tiap \unit{kW.h}.
\end{exercise}

\begin{exercise}[$\star$]\label{exo:g10:energy-conservation:3}
Ubahlah: \qty{1}{kW.h} ke \omterm{def:g10:energy-conservation:energy}{joule}; baterai telepon \qty{12}{W.h} ke
kilojoule; \qty{45}{MJ} (satu \omterm{def:g10:orders-of-magnitude:unit}{kilogram} bensin) ke \omterm{def:g10:energy-conservation:kwh}{kilowatt-jam};
makanan harianmu \qty{10}{MJ} ke \omterm{def:g10:energy-conservation:kwh}{kilowatt-jam}.
\end{exercise}

\begin{exercise}[$\star$]\label{exo:g10:energy-conservation:4}
Sebuah ketel \qty{2200}{W} menyala selama \qty{150}{s}. Hitunglah
\omterm{def:g10:energy-conservation:energy}{energi} yang dipakainya, dalam \omterm{def:g10:energy-conservation:energy}{joule} dan dalam \omterm{def:g10:energy-conservation:kwh}{kilowatt-jam}. Lalu
hitunglah \omterm{def:g10:energy-conservation:power}{daya} rata-rata seorang manusia yang berjalan dengan
\qty{10}{MJ} tiap hari.
\end{exercise}

\begin{exercise}[$\star$]\label{exo:g10:energy-conservation:5}
Tulislah \omterm{def:g10:energy-conservation:transfer}{rantai pengubahan}
(\cref{met:g10:energy-conservation:chain}) untuk: dinamo sepeda yang
menyalakan lampu; pemanas air bergas; panel surya yang mengisi telepon.
Tandailah kebocoran kalornya.
\end{exercise}

\begin{exercise}[$\star\star$]\label{exo:g10:energy-conservation:6}
Sebuah mesin membakar \qty{2.0}{kg} bensin (\qty{45}{MJ/kg}) dan
menyerahkan \qty{27}{MJ} gerak. Hitunglah \omterm{def:g10:energy-conservation:efficiency}{efisiensinya} dan \omterm{def:g10:energy-conservation:energy}{energi} yang
dilepaskan sebagai kalor. Ke mana kalor itu pergi?
\end{exercise}

\begin{exercise}[$\star\star$]\label{exo:g10:energy-conservation:7}
Seorang peloncat tebing bermassa \qty{70}{kg} melangkah dari tebing
setinggi \qty{20}{m} (hambatan udaranya diabaikan). Dengan \omterm{thm:g10:energy-conservation:conservation}{kekekalan
energi}, hitunglah \omterm{def:g10:energy-conservation:kinetic}{energi kinetiknya} saat mencapai air dan laju
masuknya, dalam \unit{m/s} dan \unit{km/h}.
\end{exercise}

\begin{exercise}[$\star\star$]\label{exo:g10:energy-conservation:8}
Sebuah meteran menunjukkan \qty{41236}{kW.h} pada 1 Maret dan
\qty{41562}{kW.h} 30 hari kemudian: berapa \omterm{def:g10:energy-conservation:energy}{energi} yang dipakai, berapa
rata-rata hariannya, berapa tagihannya pada $0.26$ euro tiap
\unit{kW.h} ditambah $12$ euro tetap, dan berapa \omterm{def:g10:energy-conservation:power}{daya} rata-ratanya
dalam \omterm{def:g10:energy-conservation:power}{watt}?
\end{exercise}

\begin{exercise}[$\star\star$]\label{exo:g10:energy-conservation:9}
Sebuah bohlam tua \qty{60}{W} hanya mengubah $5\%$ konsumsinya menjadi
cahaya; sebuah LED memberikan cahaya yang sama dari \qty{9}{W}.
Hitunglah \omterm{def:g10:energy-conservation:power}{daya} cahaya bohlam itu dan \omterm{def:g10:energy-conservation:efficiency}{efisiensi} LED-nya; lalu, pada
\qty{4}{h} tiap hari, konsumsi tahunan masing-masing dan penghematannya
pada $0.25$ euro tiap \unit{kW.h}.
\end{exercise}

\begin{exercise}[$\star\star$]\label{exo:g10:energy-conservation:10}
Sebuah bendungan mengalirkan \qty{150}{m^3} air (\qty{1.5e5}{kg}) tiap
sekon melalui terjunan setinggi \qty{40}{m}. Hitunglah \omterm{def:g10:energy-conservation:energy}{energi} potensial
yang dilepaskan tiap sekon --- yaitu \omterm{def:g10:energy-conservation:power}{daya} yang tersedia --- lalu \omterm{def:g10:energy-conservation:power}{daya}
listriknya pada $\eta = 0.90$, dan berapa rumah tangga
berdaya rata-rata \qty{525}{W} yang dipasoknya.
\end{exercise}

\begin{exercise}[$\star\star$]\label{exo:g10:energy-conservation:11}
Sebuah bandul \qty{0.20}{kg} dilepaskan \qty{5.0}{cm} di atas titik
terendahnya: hitunglah lajunya di sana. Setelah beberapa menit bandul
itu tergantung diam: ke mana \omterm{def:g10:energy-conservation:energy}{energinya} pergi, dan mengapa hal itu tidak
bertentangan dengan
\cref{thm:g10:energy-conservation:conservation}?
\end{exercise}

\begin{exercise}[$\star\star\star$]\label{exo:g10:energy-conservation:12}
Sebuah pembangkit batu bara ($\eta_1 = 0.38$), jaringannya
($\eta_2 = 0.94$) dan sebuah pengisi \omterm{def:g10:energy-conservation:power}{daya} ($\eta_3 = 0.85$) mengisi
ulang baterai telepon \qty{12}{W.h}. Hitunglah \omterm{def:g10:energy-conservation:efficiency}{efisiensi}
keseluruhannya; \omterm{def:g10:energy-conservation:forms}{energi kimia} yang dikonsumsi di pembangkitnya, dalam
\unit{W.h} dan kilojoule; serta batu bara (\qty{25}{MJ/kg}) tiap
pengisian, lalu tiap tahun dengan pengisian harian.
\end{exercise}

\begin{exercise}[$\star\star\star$]\label{exo:g10:energy-conservation:13}
Sebuah sambaran petir membawa sekitar \qty{5}{GJ}. Ubahlah ke
\omterm{def:g10:energy-conservation:kwh}{kilowatt-jam} dan ke hari bagi rumah tangga berkonsumsi
\qty{12.6}{kW.h} tiap hari. Petir itu diserahkan dalam sekitar
\qty{100}{\micro s}: hitunglah \omterm{def:g10:energy-conservation:power}{dayanya}, bandingkan dengan \omterm{def:g10:energy-conservation:power}{daya} listrik
rata-rata umat manusia (\qty{3}{TW}), lalu jelaskan mengapa tidak ada
yang memanen petir.
\end{exercise}

\begin{exercise}[$\star\star\star$]\label{exo:g10:energy-conservation:14}
Sebuah mobil \qty{1300}{kg} mengerem dari \qty{110}{km/h}
(\qty{30.6}{m/s}) sampai berhenti. Hitunglah \omterm{def:g10:energy-conservation:kinetic}{energi kinetik} yang
terbuang; bentuk apa yang diambilnya, dan di mana? Sebuah mobil listrik
memungut kembali $60\%$ ke dalam baterainya: berapa \omterm{def:g10:energy-conservation:energy}{energinya} tiap
perhentian, dan berapa perhentian yang diperlukan untuk mengisi satu
\omterm{def:g10:energy-conservation:kwh}{kilowatt-jam}?
\end{exercise}

\begin{exercise}[$\star\star\star$]\label{exo:g10:energy-conservation:15}
Seorang pengendara sepeda beserta sepedanya (\qty{85}{kg} seluruhnya)
mendaki sebuah celah gunung dan memperoleh \qty{1200}{m} ketinggian;
otot mengubah makanan dengan $\eta \approx 0.25$.
Hitunglah \omterm{def:g10:energy-conservation:energy}{energi} mekanis pendakian itu, \omterm{def:g10:energy-conservation:energy}{energi} makanan yang
dituntutnya, setaranya dalam mangkuk sereal \qty{1.5}{MJ}, dan massa
bensin (\qty{45}{MJ/kg}) yang menyimpan \omterm{def:g10:energy-conservation:energy}{energi} yang sama. Berikan
ulasanmu.
\end{exercise}

\section{Soal: Menyalakan sebuah rumah selama sehari}

\begin{problem}[{Soal akhir pekan --- dari semangkuk sereal sampai
panel di atap: audit energi lengkap satu hari biasa di rumah, dan
berapa harga satu joule bergantung pada siapa yang menjualnya}]
\label{pb:g10:energy-conservation:1}
Pada suatu Selasa yang biasa, sebuah keluarga berniat mencari tahu ke
mana listriknya pergi, berbekal bab ini, meteran di lorong rumah dan
punggung lembar tagihan. Putusan malam ini: satu hari mereka setara
satu \omterm{def:g10:orders-of-magnitude:unit}{kilogram} bensin, enam belas \omterm{def:g10:orders-of-magnitude:unit}{meter} persegi cahaya matahari --- atau
lima pengendara sepeda yang mengayuh sepanjang siang dan malam.

\medskip
\noindent\textbf{Bagian I --- Auditnya.}
Daftar hari itu: lemari es \qty{100}{W} selama \qty{24}{h}; pemanas air
\qty{2400}{W} selama \qty{2.5}{h}; oven \qty{2000}{W} selama
\qty{1.0}{h}; mesin cuci \qty{500}{W} selama \qty{2.0}{h}; penerangan
\qty{60}{W} selama \qty{5.0}{h}; layar \qty{150}{W} selama \qty{6.0}{h}.
\begin{enumerate}
  \item Hitunglah \omterm{def:g10:energy-conservation:energy}{energi} tiap alatnya untuk sehari, dalam \unit{kW.h}.
  \item Jumlahkan hari itu dalam \omterm{def:g10:energy-conservation:kwh}{kilowatt-jam}, lalu dalam megajoule;
        bandingkan dengan satu \omterm{def:g10:orders-of-magnitude:unit}{kilogram} bensin (\qty{45}{MJ}).
  \item Hitunglah \omterm{def:g10:energy-conservation:power}{daya} rata-rata rumah itu sepanjang \qty{24}{h};
        berapa bagian dari sebuah ketel yang menyala itu?
  \item Alat tunggal manakah yang paling menguasai? Hitunglah kembali
        jumlah hariannya jika pemanas airnya hanya menyala \qty{1.25}{h}.
  \item Lampu siaga menarik \qty{15}{W} secara tetap: berapa \omterm{def:g10:energy-conservation:energy}{energi}
        hariannya, berapa bagiannya dari jumlah seluruhnya, dan berapa
        biaya tahunannya pada $0.25$ euro tiap \unit{kW.h}?
\end{enumerate}

\medskip
\noindent\textbf{Bagian II --- Tagihannya.}
\begin{enumerate}[resume]
  \item Dalam 30 hari meterannya bergerak dari \qty{56214}{kW.h} ke
        \qty{56604}{kW.h}: berapa konsumsinya, dan berapa rata-rata
        hariannya jika dicocokkan dengan auditnya?
  \item Hitunglah tagihannya: $0.25$ euro tiap \unit{kW.h} ditambah
        langganan tetap $14$ euro.
  \item Berapa harga satu megajoule dari jaringan listrik? Dan dari
        bensin, pada $1.80$ euro tiap liter yang menyimpan
        \qty{34}{MJ}?
  \item Semangkuk sereal \qty{1.5}{MJ} berharga sekitar $0.50$ euro:
        berapa harga megajoulenya? Urutkan ketiga penjual \omterm{def:g10:energy-conservation:energy}{joule} itu.
\end{enumerate}

\medskip
\noindent\textbf{Bagian III --- Atapnya.}
\begin{enumerate}[resume]
  \item Cahaya matahari penuh menyerahkan sekitar \qty{1000}{W} tiap
        \omterm{def:g10:orders-of-magnitude:unit}{meter} persegi; panelnya mengubahnya dengan $\eta = 0.20$:
        berapa \omterm{def:g10:energy-conservation:power}{daya} listrik \qty{1}{m^2} pada matahari penuh?
  \item Daerah itu rata-rata memperoleh setara \qty{4.0}{h} matahari
        penuh tiap hari: berapa hasil harian \qty{1}{m^2}?
  \item Berapa luas panel yang menutupi hari \qty{12.6}{kW.h} milik
        keluarga itu? Atap selatannya menyediakan \qty{20}{m^2}: apakah
        muat?
  \item Pada bulan Desember setaranya turun menjadi \qty{1.5}{h}:
        berapa hasil panelnya, berapa bagiannya dari kebutuhan, dan apa
        yang menutupi kekurangannya?
  \item Pembuatan panelnya memakan sekitar \qty{700}{kW.h} tiap \omterm{def:g10:orders-of-magnitude:unit}{meter}
        persegi: berapa waktu pengembalian \omterm{def:g10:energy-conservation:energy}{energinya}, dengan memakai
        pertanyaan 11?
\end{enumerate}

\medskip
\noindent\textbf{Bagian IV --- Semangkuk sereal.}
\begin{enumerate}[resume]
  \item Sebuah sepeda pembangkit menghasilkan \qty{100}{W} listrik:
        berapa jam mengayuh yang menghasilkan \qty{12.6}{kW.h} ---
        berapa orang yang mengayuh \qty{24}{h} tanpa henti?
  \item Otot bekerja pada $\eta \approx 0.25$: berapa \omterm{def:g10:energy-conservation:energy}{energi} makanan di
        balik \qty{12.6}{kW.h} itu, dalam megajoule dan dalam mangkuk
        sereal?
  \item Berapa biaya hari yang dikayuh itu (mangkuknya $0.50$ euro),
        dibandingkan dengan yang ditagihkan jaringan listrik untuk hari
        yang sama?
  \item Sambaran petir \qty{5}{GJ} setara berapa hari rumah tangga?
        Mengingat lamanya \qty{100}{\micro s}, mengapa petir itu tetap
        tidak dapat menyalakan rumahnya?
  \item Joule tidak pernah hilang, namun \omterm{def:g10:energy-conservation:energy}{joule} dari jaringan listrik,
        dari bensin dan dari makanan dijual dengan harga yang sangat
        berbeda: sebenarnya apa yang kamu bayar?
  \item Penutup: nyatakan putusan hari itu dalam satu kalimat --- satu
        \omterm{def:g10:orders-of-magnitude:unit}{kilogram} bensin, enam belas \omterm{def:g10:orders-of-magnitude:unit}{meter} persegi atap yang tersinari
        matahari atau lima pengendara sepeda sepanjang siang dan malam
        --- lalu sebutkan dua besaran (satu persediaan, satu aliran)
        yang tidak boleh lagi dicampuradukkan.
\end{enumerate}
\end{problem}
